\documentclass[10pt,a4paper]{amsart}
\usepackage[margin=19mm]{geometry}
\usepackage{amsmath,amssymb,mathtools}
\usepackage[scr=rsfs,cal=euler]{mathalfa}
\usepackage{xfrac}
\usepackage[unicode,hidelinks,bookmarks=false]{hyperref}
\newcommand{\ie}{i.e.\ }
\newcommand{\cf}{cf.\ }
\newcommand{\resp}{respectively\ }
\newcommand{\Subsection}{\subsection}
\providecommand{\nobreakdash}{\nobreak\hbox{-}}
\setlength{\emergencystretch}{2em}
\multlinegap0pt
\usepackage{bm}
\usepackage{bbm}
\usepackage{dsfont}
\usepackage{xspace}
\usepackage{cancel}
\usepackage{mathtools}
\usepackage{amsthm}
\makeatletter
\@ifundefined{thm}{\newtheorem{thm}{Theorem}}{}
\makeatother
\DeclareMathOperator*{\argmax}{arg\,max}
\DeclareMathOperator{\Tr}{Tr}
\title[Matrix-Driven Identification and Reconstruction of LLM Weigh]{Matrix-Driven Identification and Reconstruction of LLM Weight Homology}
\author{Ruichong Zhang, Daniel Goldstein}
\date{Source: arXiv:2508.06309v3 (CC BY 4.0)}
\begin{document}
\maketitle

\noindent\emph{Source and licence.} Matrix-Driven Identification and Reconstruction of LLM Weight Homology. Version arXiv:2508.06309v3 (CC BY 4.0), distributed under the Creative Commons Attribution 4.0 International licence, as recorded in the arXiv licence metadata for this exact version and shown on the abstract page https://arxiv.org/abs/2508.06309v3. Full rights notice: licenses/arxiv-2508.06309-CC-BY-4.0.txt.\par\smallskip

\noindent\textit{Editorial scope.} Counted unit: the Thm whose source label is \texttt{None} in the article (source0.tex lines 1330--1343, source label \texttt{None}), delivered together with the article's own complete original proof of it (source0.tex lines 1344--1378, one \texttt{proof} environment of 35 lines). No mathematics is written, repaired or re-derived: statement and proof are the article's own text line for line. Required context retained uncounted: none. Every definition, display and label of those lines is retained unchanged. Omitted from this package and not redistributed: 1--1329, 1379--1885. Nothing inside a retained excerpt is omitted or altered: every retained body is delivered exactly as the article's lines are written, so no omission receipt of any kind accompanies this package. Citations: the retained excerpts carry no citation command at all, so no bibliography entry of the article is transcribed and no reference is redistributed. Reference adaptation: none was needed and none was made; every reference inside the delivered unit resolves inside it, so no unresolved reference or citation is produced. Printed slips kept literal: no printed word, symbol, hypothesis or number of the counted statement or of its proof was altered. Layout and class: the article's own class is \texttt{article} with packages including icml2026, amsmath, amsfonts, bm, hyperref, url, graphicx, subfigure, algorithm, algpseudocode, amsthm, booktabs, subcaption, array; the wrapper is the lane's pinned \texttt{amsart} editorial wrapper at 10pt on A4, which restates the theorem-like environments the retained text needs and adds the article's own macro definitions that the retained text needs (\texttt{\textbackslash Tr}, \texttt{\textbackslash argmax}), with extra wrapper packages (bm, bbm, dsfont, xspace, cancel, mathtools). The delivered numbering is the unit's own. Assets: no retained span contains a figure environment, a graphics-inclusion command, a PDF-page inclusion, a table or a TikZ picture; no image file is copied into this package, the package declares no asset dependency and no image file is redistributed with it. Licence: version arXiv:2508.06309v3 of this article is distributed under the Creative Commons Attribution 4.0 International licence, as recorded in the arXiv licence metadata for this exact version and shown on the abstract page https://arxiv.org/abs/2508.06309v3. A full rights notice is delivered at licenses/arxiv-2508.06309-CC-BY-4.0.txt. Characters: every retained excerpt is pure ASCII, so the delivered engine, XeLaTeX with its default Latin Modern text font, reports no missing character.
\begin{thm}
Define the trace and the corresponding optimal permutation respectively as follows:
\[
\begin{aligned}
    T(E, E') &= \max_{P \in \mathrm{Perm}(\mathrm{n})} \Tr\left(P \cdot \mathrm{Ortho}\left(E'^{\text{T}} E\right)^{\text{T}}\right), \\
    \Pi(E, E') &= \argmax_{P \in \mathrm{Perm}(\mathrm{n})} \Tr\left(P \cdot \mathrm{Ortho}\left(E'^{\text{T}} E\right)^{\text{T}}\right).
\end{aligned}
\]
Suppose an adversary applies a simultaneous scaling and permutation transformation to $E'$, yielding $E'' = \alpha E' P'$, where $P' \in \mathrm{Perm}(\mathrm{n})$ is an arbitrary permutation matrix, and $\alpha > 0$ is a positive scaling coefficient. In this case, the following properties hold:
\[
T(E, E'') = T(E, E'), \quad \Pi(E, E'') = P'^{\text{T}} \Pi(E, E').
\]
Thus, such adversarial modifications are futile.
\end{thm}

\begin{proof}
We first analyze the trace $T(E, E'')$:
\[
\begin{aligned}
    T(E, E'') &= \max_{P \in \mathrm{Perm}(\mathrm{n})} \Tr\left(P \cdot \mathrm{Ortho}\left((E''^{\text{T}} E)^{\text{T}}\right)\right) \\
    &= \max_{P \in \mathrm{Perm}(\mathrm{n})} \Tr\left(P \cdot \mathrm{Ortho}\left(((\alpha E' P')^{\text{T}} E)^{\text{T}}\right)\right) \\
    &= \max_{P \in \mathrm{Perm}(\mathrm{n})} \Tr\left(P \cdot \mathrm{Ortho}\left(\alpha E^{\text{T}} E' P'\right)\right).
\end{aligned}
\]
Since scaling by $\alpha > 0$ does not affect the orthogonal factor in polar decomposition ($\mathrm{Ortho}$), we can simplify further:
\[
\begin{aligned}
    T(E, E'') &= \max_{P \in \mathrm{Perm}(\mathrm{n})} \Tr\left(P \cdot \mathrm{Ortho}\left((E^{\text{T}} E') P'\right)\right) \\
    &= \max_{P \in \mathrm{Perm}(\mathrm{n})} \Tr\left(P \cdot \mathrm{Ortho}(E^{\text{T}} E') P'\right) \\
    &= \max_{P \in \mathrm{Perm}(\mathrm{n})} \Tr\left(P' P \cdot \mathrm{Ortho}\left(E^{\text{T}} E'\right)\right).
\end{aligned}
\]
Let $R = P' P$. Since $P'$ is a fixed permutation matrix, as $P$ ranges over all permutations in $\mathrm{Perm}(\mathrm{n})$, so does $R$. Substituting $R$ into the expression, we obtain:
\[
\begin{aligned}
    T(E, E'') &= \max_{R \in \mathrm{Perm}(\mathrm{n})} \Tr\left(R \cdot \mathrm{Ortho}\left(E^{\text{T}} E'\right)\right) \\
    &= \max_{R \in \mathrm{Perm}(\mathrm{n})} \Tr\left(R \cdot \mathrm{Ortho}\left((E'^{\text{T}} E)^{\text{T}}\right)\right) \\
    &= T(E, E').
\end{aligned}
\]
Next, we examine the optimal permutation $\Pi(E, E'')$. Recall that the trace achieves its maximum when $R = \Pi(E, E')$. From the substitution $R = P' P$, it follows that:
\[
P' P = \Pi(E, E') \! \implies \! P = P'^{-1} \Pi(E, E') = P'^{\text{T}} \Pi(E, E').
\]
Thus, the optimal permutation for $E''$ is given by:
\[
\Pi(E, E'') = P'^{\text{T}} \Pi(E, E').
\]
This completes the proof.
\end{proof}

\end{document}
